\chapter{Penutup}

\section{Kesimpulan}

Penelitian ini membangun sistem monitoring detak jantung single-lead berbasis \textit{edge computing} dengan model CNN \textit{hybrid} INT8 yang berjalan \textit{on-device} pada ESP32-S3. Preprocessing sinyal ECG dan klasifikasi aritmia dijalankan di perangkat, sedangkan \textit{pipeline} MQTT dengan bufer lokal di PSRAM untuk mengirim hasil klasifikasi per detak telah dirancang. Berdasarkan perancangan, implementasi, dan eksperimen yang telah dilakukan, kesimpulan terhadap tujuan yang hasilnya dilaporkan pada tahap ini diuraikan sebagai berikut.

Preprocessing di ESP32-S3 berjalan secara \textit{real-time}. Beban rata-rata hanya 82~$\mu$s per sampel atau 3,0\% anggaran waktu tanpa sampel yang hilang, sedangkan puncak beban sesaat 33,4~ms diserap oleh ADC yang dibaca di dalam rutin interupsi (ISR) beserta antreannya. Segmentasi hasil detektor di perangkat memberi F1 0,6499 dibandingkan 0,6594 pada segmentasi anotasi. Sinyal tubuh dari AD8232 juga layak sebagai masukan, karena keempat rekaman resmi dari satu sesi perekaman lolos indeks kualitas sinyal pada seluruh jendela 10 detik. Gerbang kalibrasi berbasis kontak elektroda membuka pemantauan sekitar 60 detik setelah kontak baik dan menahan detak selama kontak buruk. Seluruh target jalur preprocessing karena itu terpenuhi.
% sumber: bab4 tab:beban-komputasi, tab:segmentasi-deteksi, tab:kualitas-rekaman (4 dari 4 Berhasil), subsec:hasil-tujuan1 butir Gerbang Kalibrasi dan Kualitas Sinyal (sekitar 60 detik, kontak buruk ditahan), tab:ketercapaian (preprocessing Tercapai)

Model INT8 pada DS2 yang berisi 22 pasien yang tidak terlihat saat pelatihan mencapai \textit{recall} 0,7105, F1 0,6604, dan AUC 0,9344. F1 tersebut 100,14\% F1 model \textit{float32}, dengan ukuran model 22,94~KB. Inferensinya di ESP32-S3 memakan 26,6~ms per detak dengan \textit{tensor arena} 12.948 byte. Karena itu target kuantisasi, latensi, dan memori terpenuhi tanpa ketergantungan pada \textit{cloud}. Kelemahan model terletak pada kelas S dengan \textit{recall} 0,4630 pada DS2.
% sumber: bab4 tab:hasil-ds2, tab:delta-int8, tab:sumber-daya, subsec:hasil-tujuan2 butir Evaluasi DS2 (recall S 0,4390) dan butir Konsumsi Daya, tab:ketercapaian

Selain ketercapaian tujuan, penelitian ini menghasilkan dua temuan teknis yang tidak terlihat dari rancangan awal.
Pertama, augmentasi \textit{jitter} yang meniru \textit{error} detektor sendiri bekerja sebagai regularisasi. Residu detektor Pan-Tompkins pada DS1 berbentuk inti yang tajam dengan ekor yang berat, berbeda dari \textit{jitter} seragam yang dipakai penelitian terdahulu. Karena itu augmentasi mengambil sampel dari residu terukur tersebut. Pada ablasi tiga \textit{seed}, F1 DS2 naik dari 0,5660 menjadi 0,6354 walaupun DS2 sendiri tidak diberi \textit{jitter}, sehingga manfaatnya bukan sekadar simulasi kondisi lapangan. Catatannya, DS2 ikut menilai ablasi ini. Kenaikan tersebut karena itu berpotensi sedikit optimistis. Implikasinya, konstanta empiris dari penelitian lain sebaiknya diukur ulang pada detektor sendiri sebelum diadopsi.
% sumber: bab4 tab:ablasi baris bersih dan jitter (0,5660, 0,6354, tiga seed), subsec:hasil-tujuan2 butir Ketahanan terhadap Error Segmentasi (residu inti tajam dan ekor berat, berbeda dari jitter seragam Dias) dan butir Ablasi (DS2 tidak diberi jitter, regularisasi), butir Evaluasi DS2 (DS2 ikut menilai ablasi, berpotensi sedikit optimistis)

Kedua, telemetri akuisisi dapat menyesatkan, sehingga kualitas sinyal tidak boleh dinilai dari keberadaan R-\textit{peak}. Detektor Pan-Tompkins memakai ambang adaptif terhadap sinyal yang ada dan tetap melaporkan R-\textit{peak} ketika biopotensial sudah hilang. Tanpa elektroda dan tanpa \textit{replay}, perangkat mengeluarkan sekitar 2 detak per detik dengan 57\% di antaranya diklasifikasi Aritmia. \textit{Flag lead-off} AD8232 pun dapat dipenuhi oleh satu jalur impedansi rendah saja dan karena itu tidak menjamin adanya sinyal. Sebaliknya, ayunan sinyal ECG yang kuat tetap tinggi ketika kontak elektroda memburuk. Penentu kualitas yang benar karena itu adalah gabungan status \textit{lead-off}, ayunan ADC absolut, dan dengung 50 Hz, dan gabungan inilah yang dipakai sebagai gerbang kualitas di perangkat.
% sumber: bab4 subsec:hasil-tujuan1 butir Akuisisi Sinyal Tubuh (perangkat tanpa elektroda sekitar 2 detak per detik, 57% Aritmia, ambang sekitar 150 count, ayunan absolut sebagai gerbang kualitas); bab2 subbab AD8232 (lead-off cukup dipenuhi satu jalur impedansi rendah); model/CLAUDE.md jebakan "Tanpa elektroda, alat mengarang aritmia\"

Secara keseluruhan, jalur preprocessing \textit{real-time} dan jalur model TinyML \textit{on-device} tercapai (Tabel~\ref{tab:ketercapaian}). Jalur \textit{pipeline edge-to-cloud} yang resilien telah dirancang, dan pengujiannya dilaporkan pada tahap berikutnya. Sistem ini berstatus \textit{proof of concept} di lingkungan non-klinis dan bukan alat diagnosis.

\section{Saran}

Berdasarkan keterbatasan sistem dan pekerjaan yang tersisa, saran pengembangan selanjutnya diuraikan pada butir berikut.
\begin{enumerate}
  \item \textbf{Validasi pada subjek yang lebih beragam.} Rekaman sinyal tubuh perlu diperluas ke lebih banyak subjek dan sesi, sebaiknya disertai alat acuan detak jantung, agar variasi sinyal antarindividu terwakili. Pada sinyal tubuh tersebut, gerbang kalibrasi berbasis kontak elektroda perlu diuji ulang dengan transisi fase yang tercatat sebagai log, termasuk ambang dengung yang dapat berbeda antarsubjek dan antarruangan.
  % sumber: bab4 subsec:analisis-keterbatasan butir 1, 2, dan 6; subsec:hasil-tujuan1 butir Gerbang Kalibrasi dan Kualitas Sinyal (satu subjek, transisi kontak baik tanpa log serial; ambang dengung 500 dari tiga sesi satu subjek)
  \item \textbf{Pengujian \textit{pipeline edge-to-cloud}.} Jalur \textit{pipeline} MQTT dengan bufer lokal di PSRAM telah dirancang, dan hasil pengujiannya belum termasuk dalam laporan tahap ini. Pengujian berikutnya menilai PDSR, median latensi \textit{end-to-end}, bufer lokal tanpa \textit{overflow}, waktu \textit{flush}, dan deduplikasi terhadap target pada Tabel~\ref{tab:pengujian-pipeline}, termasuk skenario putus koneksi dan tampilan data pada \textit{dashboard} ThingsBoard. Pengujian ini juga sebaiknya memakai sinyal tubuh, bukan hanya sinyal \textit{replay}.
  % sumber: bab3 tab:pengujian-pipeline (enam target); bab4 tab:parameter-eksperimen kategori f; hasil pipeline dibuang dari progres-pa-1 atas keputusan user 1 Nov 2026, tetap ada di laporan/buku-pa/bab/bab4.tex subsec:hasil-tujuan3
  \item \textbf{Peningkatan deteksi kelas S.} Dugaan bahwa periode refraktori 200~ms memblokir detak S yang prematur dan menekan \textit{sensitivity} detektor kelas S ke 0,6494 perlu diuji dengan mengubah periode refraktori. Pertukaran \textit{learning rate} $10^{-3}$ terhadap $10^{-4}$ pada Tabel~\ref{tab:ablasi} juga dapat ditinjau ulang bila prioritas kelas berubah. Konsekuensinya, model pada \textit{firmware} dibuat ulang serta divalidasi ulang di perangkat.
  % sumber: bab4 tab:sensitivity-detektor (S 0,6494) dan butir Segmentasi di Perangkat (dugaan periode refraktori), subsec:hasil-tujuan2 butir Ablasi (learning rate)
  \item \textbf{Perangkat \textit{wearable}.} Rakitan saat ini perlu dipindahkan ke PCB yang menyatukan seluruh modul, lalu dikemas agar perangkat dapat dipakai secara portabel. Klasifikasi juga dapat diperluas dari biner menjadi penamaan jenis aritmia bila aturan penamaannya dapat dinilai pada data validasi.
  % sumber: hardware/mainboard (PCB satu papan dengan catu baterai dan INA219); firmware/README.md bagian Rangkaian catu baterai; model/docs/README.md "Eksperimen yang tidak diadopsi" (aturan penamaan V/S tidak lolos sertifikasi VAL)
\end{enumerate}
